% 
% Updated: May 27, 1997 
% 
% \documentclass[12pt]{amsart}
\documentstyle[12pt]{article}
\font\tit=cmbx12
\font\bigbf=cmbx12 scaled \magstep1
\parindent=0pt
\setlength{\textwidth}{160mm}
\setlength{\textheight}{200mm}
\addtolength{\headheight}{1.5pt}
\setlength{\oddsidemargin}{2mm}
\setlength{\evensidemargin}{\oddsidemargin}
\addtolength{\textheight}{3.2cm}
\addtolength{\topmargin}{-0.7cm}
\hyphenation{Sprin-ger}

\def\R{{\rm I\kern-.20em R}}
\begin{document}
\begin{center}
{\Large \bf Presentation in IMACS '97, Maui, Hawaii}
\vskip 20pt
{\Large \bf Conference on Applications of Computer Algebra}
\vskip 20pt
{\Large \bf Session: Algorithmic Analysis of PDEs}
\vskip 20pt
{\bf \"Unal G\"okta\c{s} and Willy Hereman} \\
Department of Mathematical and Computer Sciences \\
Colorado School of Mines \\
Golden, CO 80401-1887, USA
\vskip 20pt 
{\bf Speaker:} \"Unal G\"okta\c{s}, ugoktas@lie.mines.edu
\end{center}
\vskip 3pt
\begin{center}
{\large \bf
Symbolic Computation of Conserved Densities for Nonlinear
Evolution and Differential-Difference Equations 
} 
\end{center}
A straightforward method for the computation of polynomial-type conserved
densities for systems of nonlinear evolution equations and its extension
towards differential-difference equations will be presented.
The technique is based on the scaling symmetries of the equations,
and is completely algorithmic.
The method can be implemented in any computer algebra system.

Our programs in {\it Mathematica\/} allow one to automatically
compute conserved densities.
Applied to systems with parameters, they can be used to determine
the conditions on these parameters so that a sequence of conservation
laws will exist, hence predicting integrability. 
\vskip 5pt
{\bf Software Demonstration:}
\vskip 5pt
CONDENS.M: Symbolic Computation of Conservation Laws for Systems
of Nonlinear Evolution Equations with {\it Mathematica}
\vskip 5pt
DIFFDENS.M: Symbolic Computation of Conservation Laws for Systems
of Nonlinear Differential-Difference Equations with {\it Mathematica}
\vskip 5pt
The programs {\it condens.m} and {\it diffdens.m} will be demonstrated with
several examples, including equations with parameters, for which
the compatibility conditions for the existence of conserved densities
will be computed.
\vfill
\end{document}
\end
\bye
